\documentclass[12pt]{book}
\usepackage{tikz}
\usepackage[paperwidth=12cm, paperheight=12cm, vmargin=0cm]{geometry}
\setlength{\parindent}{0cm}
\setlength{\parskip}{12pt}
\renewcommand{\familydefault}{cmss}

\newcommand{\imgpage}[1]{
    \newpage
    \tikz[remember picture,overlay]
    \node[inner sep=0pt] at (current page.center){
    \includegraphics[width=12cm,height=12cm]{#1}
    };
    \newpage
}

\title{Madame Je-Veux-Pas-Dormir}
\author{Guillaume Noël-Martel}
\date{\today}

\begin{document}
\large
\pagenumbering{gobble}

\hspace{0pt}\vfill

Madame Je-Veux-Pas-Dormir se planifiait toujours de grosses journées. Ce mardi ne serait pas une
exception.

Dès l'aurore, elle se réveilla, bondit de son lit, et commença immédiatement sa journée.
Chaque minute passée à dormir, c'est une minute gâchée qu'elle disait toujours.
\hspace{0pt}\vfill

\imgpage{img/page-1.jpg}

\hspace{0pt}\vfill

En ce mardi, elle débuta sa journée comme tant d'autres:

Après avoir fait sa toilette, elle s'assit devant un bon déjeuner de crêpes et de confiture.

À la fin de son déjeuner, elle bâilla. Puis alors qu'elle pensait en être sortie, elle bâilla
encore.

Ohho. Ça veut dire qu'il est temps d'aller faire de la musique!

\hspace{0pt}\vfill

\imgpage{img/page-2.jpg}

\hspace{0pt}\vfill

Vite, car il ne fallait pas que le sommeil la rattrape, madame Je-Veux-Pas-Dormir quitta son
déjeuner.

Elle s'assit devant son piano. Un peu de musique et de danse, ça c'est une bonne idée.

Oui, une petite balade, pas trop vite, et puis...

\hspace{0pt}\vfill

\imgpage{img/page-3.jpg}

\hspace{0pt}\vfill

Madame Je-Veux-Pas-Dormir ouvrit les yeux.

Tiens, c'est drôle elle est allongée sur son sofa. Elle n'était pas en train de faire du piano?

...

Ohho. Elle s'était encore fait avoir.

Il était déjà l'heure du dîner!

\hspace{0pt}\vfill

\imgpage{img/page-4.jpg}

\hspace{0pt}\vfill

Plus tard cette journée-là, Madame Je-Veux-Pas-Dormir décida de faire une bonne marche dehors.

Sur le chemin, elle croisa son bon ami, monsieur Je-Suis-Pas-Fatigué.

- Vous avez l'air de bien triste mine, lui dit monsieur Je-Suis-Pas-Fatigué.

- Hélas, lui dit madame Je-Veux-Pas-Dormir. Je dois dormir tout les jours, alors je n'ai jamais le
temps de tout faire.

- Vous connaissez le café? Lui répondit monsieur Je-Suis-Pas-Fatigué. Peut-être que ça vous
permettrait d'arrêter de dormir?

\hspace{0pt}\vfill

\imgpage{img/page-5.jpg}

\hspace{0pt}\vfill

Le soir, arrivé 19h00, c'est comme tout les jours, madame Je-Veux-Pas-Dormir se met à bâiller et à
fermer des yeux.

Mais cette fois-ci, elle en a assez. Elle se rappelle la recommendation de monsieur
Je-Suis-Pas-Fatigué.

Elle a une idée! Ce serait la fin de toute ce sommeil néfaste.

\hspace{0pt}\vfill

\imgpage{img/page-6.jpg}

\hspace{0pt}\vfill

Madame Je-Veux-Pas-Dormir appelle monsieur Je-Suis-Pas-Fatigué.

- Bonsoir, annonce madame Je-Veux-Pas-Dormir. Je vous invite chez moi! Apportez le café, et ensemble
nous pourront passer toute la nuit ensembles. Pas besoin de dormir!

- Oh mais c'est une excellente idée, lui répond monsieur Je-Suis-Pas-Fatigué. J'arrive dans quelques
minutes.

\hspace{0pt}\vfill

\imgpage{img/page-7.jpg}

\hspace{0pt}\vfill

Une fois arrivé, monsieur Je-Suis-Pas-Fatigué prépare vite le café. Ils se versent tout deux de
grandes tasses, et boivent goulûment. Une grande nuit se prépare.

\hspace{0pt}\vfill

\imgpage{img/page-8.jpg}

\hspace{0pt}\vfill

Toute la nuit, madame Je-Veux-pas-Dormir et monsieur Je-Suis-Pas-Fatigué peuvent faire toutes les
choses qu'ils veulent.

Madame Je-Veux-Pas-Dormir propose d'abord de faire un grand casse-tête de cent pièces. Pour une fois
elle aura le temps de le faire au complet d'un coup.

\hspace{0pt}\vfill

\imgpage{img/page-9.jpg}

\hspace{0pt}\vfill

Le long des activités, une bonne nuit remplie s'écoule. ALors que le ciel commence un peu à
s'éclairer, monsieur Je-Suis-Pas-Fatigué prépare de nouveau du café.

Madame Je-Veux-Pas-Dormir est épatée! Son plus grand problème est définitivement réglé. Plus besoin
de dormir.

- Bonne journée annonce-t-elle à monsieur Je-Suis-Pas-Dormir tandis qu'il repart chez lui.

\hspace{0pt}\vfill

\imgpage{img/page-10.jpg}

\hspace{0pt}\vfill

Madame Je-Veux-Pas-Dormir commence donc sa nouvelle journée. Elle prépare ses crêpes pour son
déjeuner. Ohho. Elle échappe ses crêpes! Ça ne lui arrive pas d'habitude pourtant.

Elle s'assit ensuite pour manger. N'oublie-t-elle pas quelque chose?

\hspace{0pt}\vfill

\imgpage{img/page-11.jpg}

\hspace{0pt}\vfill

Après un déjeuner un peu triste, madame Je-Veux-Pas-Dormir ne ressent toujours pas de fatigue! Le
café fonctionne à merveille.

Elle en profite donc et va jouer un peu de piano.

Do. Do. Do. Do... Do?

Mais elle n'a idée quoi jouer...

\hspace{0pt}\vfill

\imgpage{img/page-12.jpg}

\hspace{0pt}\vfill

La journée s'écoule drôlement. Madame Je-Veux-Pas-Dormir a plus de temps, mais elle a l'impression
que rien ne va correctement. Étrange.

Arrivé au soir, elle rappelle monsieur Je-Suis-Pas-Fatigué.

Ça sonne.

Ça sonne encore.

Monsieur Je-Suis-Pas-Fatigué ne répond pas.

- Bon, on se reprendra une autre fois se dit donc madame Je-Veux-Pas-Dormir.

\hspace{0pt}\vfill

\imgpage{img/page-13.jpg}

\hspace{0pt}\vfill

Madame Je-Veux-Pas-Dormir commence donc sa nuit: Elle se fait une grande tasse de café, qu'elle boit
goulûment.

Après la première tasse, elle se sent un peu fébrile, alors elle en boit une deuxième.

\hspace{0pt}\vfill

\imgpage{img/page-14.jpg}

\hspace{0pt}\vfill

Madame Je-Veux-Pas-Dormir s'assoit ensuite pour faire un nouveau grand casse-tête de cent pièces.

Elle n'est pas sûr de ce qui se passe, mais ça ne se passe pas bien, mais alors pas bien du tout!

Elle a du mal à bien distinguer où vont les pièces, elle a de la misère à les mettre en place, et
ses yeux sont lourds.

Si lourds.

\hspace{0pt}\vfill

\imgpage{img/page-15.jpg}

\hspace{0pt}\vfill

Madame Je-Veux-Pas-Dormir, s'installe sur son sofa. Elle se dit qu'elle va reposer ses yeux, pas
trop longtemps. Son pendule dit qu'il est tout juste 20h00.

Elle ferme un oeil, puis elle ferme l'autre, puis dans cinq minutes elle pourra les rouvrir.

Cinq minutes?

Lorsque madame Je-Veux-Pas-Dormir, elle regarde son pendule, elle se demande si le temps a reculé.
Vraiment?

\hspace{0pt}\vfill

\imgpage{img/page-16.jpg}

\hspace{0pt}\vfill

Le pendule dit 7h55!

\hspace{0pt}\vfill

\end{document}
